\documentclass[9pt]{extarticle}
\usepackage[margin=0.65in, headheight=14pt]{geometry}
\usepackage{enumitem}
\usepackage{listings}
\usepackage{hyperref}
\usepackage{xcolor}
\usepackage{fancyhdr}
\usepackage{titlesec}
\usepackage{tikz}
\usepackage{forest}
\usetikzlibrary{arrows.meta, positioning, fit, backgrounds, calc}

\hypersetup{colorlinks=true, linkcolor=blue, urlcolor=blue}
\setlength{\parindent}{0pt}
\setlength{\parskip}{3pt}
\titlespacing*{\section}{0pt}{8pt}{3pt}
\titlespacing*{\subsection}{1.5em}{6pt}{2pt}
\setlist[itemize]{leftmargin=3em}
\setlist[enumerate]{leftmargin=3em}

\title{\vspace{-1.2cm}\huge\textbf{Getting Started}\vspace{-0.5em}}
\author{\normalsize Autobook}
\date{\vspace{-1.5em}}

\pagestyle{fancy}
\fancyhf{}
\fancyhead[L]{\small Getting Started}
\fancyhead[R]{\small Autobook}
\fancyfoot[C]{\small \thepage}
\renewcommand{\headrulewidth}{0.4pt}
\renewcommand{\footrulewidth}{0pt}

\lstset{
  basicstyle=\ttfamily\small,
  backgroundcolor=\color{gray!10},
  frame=single,
  framerule=0pt,
  breaklines=true,
  xleftmargin=3em,
  xrightmargin=8pt,
  aboveskip=4pt,
  belowskip=4pt,
}

\newcommand{\cmd}[1]{\colorbox{gray!10}{\texttt{#1}}}

\begin{document}
\maketitle
\thispagestyle{fancy}

One-time setup, how the pipeline works, and the nanochat fork workflow.

\leftskip=0pt
\section*{1. Prerequisites and One-Time Setup}
\leftskip=1.5em

Three accounts needed: GitHub (you already have this), Modal (\url{https://modal.com}) for GPU compute, and Weights \& Biases (\url{https://wandb.ai}) for experiment tracking. Create your Modal and W\&B accounts, then send your W\&B username to a teammate so they can add you to project \texttt{490-autobook-a3}. Modal gives \$30 free credits on signup; we also have a shared \$500 academic credit pool.

\subsection*{1.1 Install W\&B and get your API key}
\leftskip=3em

Run \cmd{pip install wandb}. Find your API key at \url{https://wandb.ai/authorize} --- you will need it in the next step. You also need the local package to fetch training metrics (loss, bpb) for your writeup.

\subsection*{1.2 Install Modal CLI and authenticate}
\leftskip=3em

Run \cmd{pip install modal}, then \cmd{modal token new} (opens browser to authenticate). Create a secret so Modal can log to W\&B on your behalf: \cmd{modal secret create wandb-secret WANDB\_API\_KEY=<your-key>}. Verify with \cmd{modal secret list} --- should show \cmd{wandb-secret}.

\leftskip=0pt
\section*{2. How the Pipeline Works}
\leftskip=1.5em

You write a YAML config. The pipeline reads it, sends training and eval jobs to Modal GPUs, and writes results to a Modal Volume and W\&B. You never need to edit the pipeline code --- everything is controlled via your config.

\subsection*{2.1 Files and call graph}
\leftskip=3em

\begin{forest}
  for tree={
    font=\ttfamily\small,
    grow'=0,
    child anchor=west,
    parent anchor=south,
    anchor=west,
    calign=first,
    edge path={
      \noexpand\path [draw, -{|}] (!u.south west) +(7.5pt,0) |- node[fill,inner sep=1.25pt] {} (.child anchor)\forestoption{edge label};
    },
    before typesetting nodes={
      if n=1 {insert before={[,phantom]}} {}
    },
    fit=band,
    before computing xy={l=15pt},
    s sep=4pt,
  }
  [a3/shared/scripts/
    [main.py \normalfont{\textit{--- parses YAML, calls} \texttt{run.remote()}}]
    [shared/
      [infra.py \normalfont{\textit{--- app, image, volume, secrets, setup()}}]
      [helpers.py \normalfont{\textit{--- pure Python utilities}}]
    ]
    [runners/
      [run.py \normalfont{\textit{--- orchestrator: spawns stages then evals}}]
      [train.py \normalfont{\textit{--- one training stage per GPU}}]
      [evaluate.py \normalfont{\textit{--- one checkpoint eval per GPU}}]
    ]
  ]
\end{forest}

\vspace{6pt}
\begin{center}
% --- Custom palette ---
\definecolor{powderblush}{HTML}{E8C4C4}
\definecolor{parchment}{HTML}{F1E9D2}
\definecolor{silver}{HTML}{C0C0C0}
\definecolor{greyolive}{HTML}{A2A38B}
\definecolor{taupe}{HTML}{8B7E74}
\definecolor{sunflowergold}{HTML}{FFD700}
\definecolor{darkteal}{HTML}{005F5F}
\begin{tikzpicture}[
  node distance=5mm and 6mm,
  base/.style={draw=taupe, rounded corners=2pt, minimum height=6mm, font=\small\ttfamily},
  pyfile/.style={base, fill=powderblush},
  worker/.style={base, fill=greyolive!30, minimum width=14mm},
  wbnode/.style={base, fill=sunflowergold!50, minimum width=18mm},
  volnode/.style={base, fill=darkteal!25, minimum width=18mm},
  arr/.style={-{Stealth[length=2.5mm]}, thick, color=taupe},
  dataarr/.style={-{Stealth[length=2mm]}, thick, color=taupe, opacity=0.3},
  lbl/.style={font=\sffamily\scriptsize\itshape, text=black!60},
]
  % === LOCAL ===
  \node[pyfile] (yaml) {your.yaml};
  \node[pyfile, right=8mm of yaml] (main) {main.py};
  \draw[arr] (yaml) -- (main);
  \node[lbl, above=1pt of main, text=white] {parses YAML};

  % === MODAL: orchestrator ===
  \node[pyfile, right=30mm of main, minimum width=16mm] (run) {run.py};
  \draw[arr] (main) -- node[lbl, above] {.remote()} (run);
  \node[lbl, above=1pt of run] {orchestrator};

  % W&B — same y as run.py
  \node[wbnode, right=60mm of run] (wandb) {W\&B};
  \node[lbl, below=1pt of wandb] {\scriptsize curves};

  % Setup
  \node[worker, below=8mm of run] (setup) {setup};
  \draw[arr] (run) -- (setup);
  \node[lbl, right=2mm of setup] {tokenizer + data};

  % train.py — closer to setup's x
  \node[pyfile, minimum width=20mm, right=4.5mm of setup, yshift=-12mm] (trainpy) {train.py};

  % evaluate.py — same x as train.py, more distance below
  \node[pyfile, minimum width=20mm, below=22mm of trainpy] (evalpy) {evaluate.py};

  % Straight vertical spine from setup, horizontal branches to .py nodes
  \draw[color=taupe, thick] (setup.south) -- (setup.south |- evalpy.west);
  \draw[arr] (setup.south |- trainpy.west) -- (trainpy.west);
  \draw[arr] (setup.south |- evalpy.west) -- (evalpy.west);

  % Trains — stacked to the right of train.py
  \node[worker, right=8mm of trainpy, yshift=8mm] (t1) {train 1};
  \node[worker, below=2mm of t1] (t2) {train 2};
  \node[worker, below=2mm of t2] (t3) {train 3};
  % Arrows from train.py branching to all three trains
  \coordinate (tfork) at ([xshift=4mm]trainpy.east);
  \draw[color=taupe, thick] (trainpy.east) -- (tfork);
  \draw[arr] (tfork) |- (t1.west);
  \draw[arr] (tfork) -- (t2.west);
  \draw[arr] (tfork) |- (t3.west);

  % Evals — same x as trains, stacked below evaluate.py
  \coordinate (eval-anchor) at ([yshift=8mm]evalpy.east);
  \node[worker, anchor=center] at (t1 |- eval-anchor) (e1) {eval 1};
  \node[worker, below=2mm of e1] (e2) {eval 2};
  \node[worker, below=2mm of e2] (e3) {eval 3};

  % Arrows from evaluate.py branching to all three evals
  \coordinate (efork) at ([xshift=4mm]evalpy.east);
  \draw[color=taupe, thick] (evalpy.east) -- (efork);
  \draw[arr] (efork) |- (e1.west);
  \draw[arr] (efork) -- (e2.west);
  \draw[arr] (efork) |- (e3.west);

  % Volume — right side, between train and eval groups
  \coordinate (mid) at ($(t3.south)!0.5!(e1.north)$);
  \node[volnode] at (wandb |- mid) (vol) {Volume};
  \node[lbl, below=1pt of vol] {\scriptsize checkpoints, CSVs};

  % Low-opacity curvy: trains → W&B
  \draw[dataarr] (t1.east) to[out=15, in=170] (wandb.west);
  \draw[dataarr] (t2.east) to[out=10, in=180] (wandb.west);
  \draw[dataarr] (t3.east) to[out=5, in=190] (wandb.west);

  % Low-opacity curvy: trains → Volume
  \draw[dataarr] (t1.east) to[out=-5, in=150] (vol.west);
  \draw[dataarr] (t2.east) to[out=-10, in=160] (vol.west);
  \draw[dataarr] (t3.east) to[out=-15, in=170] (vol.west);

  % Low-opacity curvy: evals → Volume
  \draw[dataarr] (e1.east) to[out=15, in=190] (vol.west);
  \draw[dataarr] (e2.east) to[out=10, in=200] (vol.west);
  \draw[dataarr] (e3.east) to[out=5, in=210] (vol.west);

  % === BACKGROUND BOXES ===
  % Local — parchment
  \begin{scope}[on background layer]
    \definecolor{localbox}{HTML}{1B1B1B}
    \node[draw=localbox, rounded corners=4pt, fill=localbox,
          fit=(yaml)(main),
          inner sep=5mm, label={[font=\scriptsize\color{parchment!40!black}]above:local}] {};
  \end{scope}

  % Modal — silver
  \begin{scope}[on background layer]
    \node[draw=silver!60!black, dashed, rounded corners=4pt, fill=silver!20,
          fit=(run)(setup)(trainpy)(t1)(t3)(evalpy)(e1)(e3)(wandb)(vol),
          inner sep=5mm, label={[xshift=13mm, anchor=south east, font=\scriptsize\color{silver!30!black}]above right:Modal cloud}] {};
  \end{scope}
\end{tikzpicture}
\end{center}

\subsection*{2.2 Workflow}
\leftskip=3em

From your perspective, using the pipeline is two steps:

\begin{enumerate}[nosep]
  \item \textbf{Write a YAML config} describing what to train and what to evaluate. The config specifies your model depth, sequence length, GPU type, W\&B run name, and which evals to run. See the \textit{Configuring Your Experiment} doc for the full reference and templates.
  \item \textbf{Run one command}, passing your config as a parameter:
\end{enumerate}

\begin{lstlisting}
modal run --detach a3/shared/scripts/main.py --config a3/<part>/configs/<your-config>.yaml
\end{lstlisting}

That's it. Modal spins up GPU containers, trains your model, runs evals, and writes results to a Modal Volume (checkpoints, CSVs, JSONs) and W\&B (training curves). You can close your laptop --- \texttt{--detach} keeps the pipeline running on Modal. Check progress on the Modal dashboard or W\&B.

\leftskip=0pt
\section*{3. Nanochat Fork}
\leftskip=1.5em

Nanochat is an LLM training codebase by Karpathy. We maintain a fork at \url{https://github.com/seonghyunban/nanochat} so that teammates can modify the architecture without touching the original code. The pipeline clones this fork into every Modal container and checks out whichever branch or tag your config specifies.

\subsection*{How to modify vanilla nanochat}
\leftskip=3em

\begin{enumerate}[nosep]
  \item Clone the fork: \texttt{git clone https://github.com/seonghyunban/nanochat.git}
  \item Create a branch: \texttt{git checkout -b p2}
  \item Make your architecture changes, commit, and push: \texttt{git push origin p2}
\end{enumerate}

Never push to \texttt{master} --- it is frozen as the unmodified baseline.

\subsection*{How to use your modified nanochat}
\leftskip=3em

Set \texttt{nanochat\_ref} in your YAML config to your branch name:

\begin{lstlisting}
nanochat_ref: p2
\end{lstlisting}

\textbf{Push before you run.} Modal runs \texttt{git fetch} + \texttt{checkout} on the remote fork --- it does not see your local changes. Forgetting to push means Modal silently runs old code.

\subsection*{How to use vanilla nanochat}
\leftskip=3em

Use the immutable baseline tag:

\begin{lstlisting}
nanochat_ref: baseline-v0
\end{lstlisting}

This tag points to unmodified nanochat and never moves, even if someone pushes to \texttt{master}.

\end{document}
